\section{Model Transmisi Data}
\subsection{Pengertian Sinyal Dalam Komunikasi Data}
Representasi fisik dari data yang dapat merambat melalui media transmisi, baik berupa listrik, cahaya, maupun gelombang elektromagnetik.


Ada 2 jenis utama yaitu:
\begin{enumerate}
\item \textbf{Sinyal Analog} \\
  Nilainya berubah secara kontinu (bersambung/tidak terputus-putus) mengikuti besaran fisik aslinya. Sinyal analog dapat memiliki nilai berapapun dalam suatu rentang tertentu, sehingga bentuk gelombangnya tampak halus dan berkesinambungan.
\item \textbf{Sinyal Digital} \\
  Sinyal yang nilainya berupa data diskrit (terputus-putus) umumnya hanya terdiri atas 2 kondisi (0 dan 1) yang direpresentasikan sebagai 2 level tegangan berbeda, misalnya ``Tinggi'' dan ``Rendah''.
\end{enumerate}

\subsection{Sinyal ADC (Analog to Digital Converter)}
ADC adalah proses/komponen yang berfungsi mengubah sinyal analog menjadi sinyal digital, agar data tersebut dapat diproses, disimpan,, atau dikirimkan oleh perangkat digital seperti komputer.

\subsection{DAC (Digital to Analog Converter)}
DAC adalah proses/komponen yang berfungsi mengubah kembali sinyal digital menjadi sinyal analog, agar data tersebut dapat dipahami/dirasakan oleh manusia melalui perangkat output analog.

\subsection{Perbandingan Antara Transmisi Analog dan Digital}
\begin{table}[H]
  \centering
  \begin{tabular}{lll}
    \toprule
    \textbf{Parameter Perbandingan} & \textbf{Transmisi Analog} & \textbf{Transmisi Digital} \\
    \midrule
    Bentuk Sinyal & Gelombang kontinu (terus-menerus) & Sinyal diskrit (biner 0 dan 1) \\
    Ketahanan terhadap Noise & Rendah (mudah terpengaruh gangguan) & Tinggi (lebih tahan dan bisa diperbarui) \\
    Keandalan \& Kualitas & Rentan atenuasi dan distorsi jarak jauh & Kualitas tetap jernih dan stabil \\
    Penggunaan Bandwidth & Lebih hemat bandwidth & Membutuhkan bandwidth yang lebih besar \\
    Fleksibilitas \& Pemrosesan & Sulit dimodifikasi dan disimpan & Sangat mudah diproses oleh komputer \\
    Contoh Penggunaan & Radio FM/AM, TV analog lama & Komputer, internet, TV digital \\
    \bottomrule
  \end{tabular}
\end{table}
